% !Mode:: "TeX:UTF-8"
\documentclass{ctexbeamer}


\usepackage{amsmath,amssymb,amsthm}             % AMS Math
\usepackage{graphicx}
\usepackage{epstopdf}
\usepackage{tikz}
\linespread{1.3}
\usepackage{amsfonts}


\usepackage{mathrsfs}  %花写字母
 

%%%=== theme ===%%%
% \usetheme{Warsaw}
%\usetheme{Copenhagen}
%\usetheme{Singapore}
\usetheme{Madrid}
%\usefonttheme{professionalfonts}
%\usefonttheme{serif}
% \usefonttheme{structureitalicserif}
%%\useinnertheme{rounded}
%%\useinnertheme{inmargin}
\useinnertheme{circles}
%\useoutertheme{miniframes}
\setbeamertemplate{navigation symbols}{}
%\setbeamertemplate{footline}[page number]
\setbeamertemplate{footline}[frame number] 


\titlegraphic{\includegraphics[width=2cm]{tjnu.jpg}} 


%\usepackage[fontset=mac]{ctex}
%\usepackage{ctex}
\usepackage{xcolor}
\newcommand{\red}[1]{\textcolor{red}{#1}}
\newcommand{\blue}[1]{\textcolor{blue}{#1}}
\newcommand{\green}[1]{\textcolor{green}{#1}}



\setbeamertemplate{theorems}[numbered]
\newtheorem{thm}{定理}
\newtheorem{lem}{引理}
\newtheorem{exa}{例}
\newtheorem*{theo}{定理}
\newtheorem*{conj}{猜想}
\newtheorem*{defi}{定义}
\newtheorem*{coro}{推论}
\newtheorem*{ex}{练习}
\newtheorem*{rem}{注}
\newtheorem*{prop}{性质}
\newtheorem*{qst}{问题}

\def\qed{\nopagebreak\hfill{\rule{4pt}{7pt}}\medbreak}
\def\pf{{\bf 证明~~ }}
\def\sol{{\bf 解~~ }}



\def\R{\mathbb{R}}
\def\Rn{\mathbb{R}^n}
\def\A{\mathscr{A}}
\def\B{\mathscr{B}}
\def\D{\mathscr{D}}
\def\E{\mathscr{E}}
\def\O{\mathscr{O}}

\def\rank{\operatorname{rank}}
\def\dim{\operatorname{dim}}
\def\0{\mathbf{0}}
\def\a{\alpha}
\def\b{\beta}
\def\r{\gamma}


\begin{document}



\title[]{论文写作指导}
\author[]{{\large 张彪} }
\institute[]{{\normalsize
		天津师范大学\\[6pt]
		zhang@tjnu.edu.cn}}

\date{}



\AtBeginSection[]
{
\begin{frame}
	\frametitle{Outline}
	\tableofcontents[currentsection]
\end{frame}
\setcounter{exa}{0}
\setcounter{equation}{0}
}



\begin{frame}
\maketitle
\end{frame}


\section{定理是什么？}


\section{优劣比较}

\begin{frame}{符号放置}
	避免以数学表达式开始一个句子，特别是如果前面的句子以数学表达式结束，否则读者难以分析句子。
	
	例如, ``$A$ is  an ill-conditioned matrix'' 
	(可能与单词``A''混淆)	
	可改为\pause
	
	 \qquad 	 \ ``The matrix $A$ is   ill-conditioned''. 
\vspace{10pt}

如果可能的话，为了同样的理由，用标点符号或者文字将数学符号隔开。

	\begin{itemize}
\item 差: If $x>1 \, f(x)<0$.\pause

\item 中: If $x>1, f(x)<0$. \pause

\item  好: If $x>1$ then $f(x)<0$.
	\end{itemize}

\vspace{10pt}

	\begin{itemize}
	\item 
差: since $p^{-1}+q^{-1}=1,\|\cdot\|_{p}$ and $\|\cdot\|_{q}$ are dual norms. \pause

\item  好: since $p^{-1}+q^{-1}=1,$ the norms $\|\cdot\|_{p}$ and $\|\cdot\|_{q}$ are dual.
	\end{itemize}
\vspace{10pt}
\end{frame}

\begin{frame}
	\begin{itemize}
	\item 
差: It suffices to show that $\|H\|_{p}=n^{1 / p}, 1 \leqslant p \leqslant 2$. \pause
	\item 
好: It suffices to show that $\|H\|_{p}=n^{1 / p}$ for $1 \leqslant p \leqslant 2$. 
\item  
好: It suffices to show that $\|H\|_{p}=n^{1 / p}(1 \leqslant p \leqslant 2)$.
	\end{itemize}
\vspace{10pt}

\begin{itemize}
	\item 
差: For $n=r(2.2)$ holds with $\delta_{r}=0$. \pause
\item 
好: For $n=r,(2.2)$ holds with $\delta_{r}=0 .$ 
\item 
好: For $n=r,$ inequality (2.2) holds with $\delta_{r}=0$.
	\end{itemize}

\end{frame}

\begin{frame}{``The''或``A''}
    
    \begin{itemize}
        \item  当一个普通名词以复数形式第一次出现时，它的前面不加the，除非它带有一限定性形容词，比如writing for sciences 和 writing for the mathematical sciences.
        \item 当一个普通名词以单数形式出现时，除了一些抽象名词外，前面要加a或the。
        比如   Let $M$ be a matroid on a ground set $E$.
        
    \end{itemize}
\end{frame}


\begin{frame}{冠词}
冠词的使用规则是 复杂的. 
{Swan} 解释了它们并认同两种最重要的规则:

\begin{itemize}
    \item  通常不要使用 $t h e$ (复数或不可数名词) 描述事情. 
    
    例: ``Mathematics is interestin''
    
    (不是 ``The mathematics is interestin'');
    
    ``Indefinite integrals do not always have closed form solutions''
    
    (不是 ``The indefinite integrals do not always have the closed form solutions'' ).
    
    \item 没有冠词就不要使用单数可数名词.
    
    例: ``the derivative is'' ``a derivative is'' , 而不是 ``derivative is''.
    
 \item 在特定环境下冠词是可选择的. 

``A matrix with the property (3.2) is well conditioned'' 和 ``A matrix with property $(3.2)$ is well conditioned'' 都是正确的.

 \item 使用冠词的错误是令人反感的, 但是它们通常不会模糊句子的含义.
\end{itemize}

\end{frame}


\begin{frame}{``The''或``A''}
在数学写作中，当宾语指代（可能）\alert{不是独一无二的}或者\alert{不存在的事物}时，使用冠词``the''是不恰当地。
改写句子，或者将冠词改为``a'', 通常能解决问题。
\vspace{10pt}

{
\begin{itemize}
\item 差: Let the Schur decomposition of $A$ be $Q T Q^{*}$. \pause
\item 好: Let a Schur decomposition of $A$ be $Q T Q^{*}$.
	\end{itemize}
}

\vspace{10pt}

\begin{itemize}
\item 差: Under what conditions does the iteration converge to the solution of $f(x)=$ 
$0 ?$  \pause
\item  好: Under what conditions does the iteration converge to a solution of $f(x)=0 ?$
\end{itemize}
\end{frame}

\begin{frame}{翻译}
\begin{itemize}
\item 
    令$a$为常数。不失一般性我们可把这个常数取成非零。

Let $a$ be a constant. Without loss of generality we may take the constant to be nonzero.
\item 
    令$A$为满秩方阵并记$A^{-1}$为$A$的逆。
    
Let $A$ be a square matrix with full rank and let $A^{-1}$ be the inverse of $A$.
\end{itemize}
\end{frame}
\end{document} 